\documentclass[11pt,a4paper]{article}
\usepackage[utf8]{inputenc}
\usepackage[T1]{fontenc}
\usepackage{lmodern,amsmath,amssymb,amsthm,mathtools,array,booktabs,enumitem,microtype}
\usepackage[margin=1in]{geometry}
\usepackage{natbib}
\usepackage{xr-hyper}
\usepackage{hyperref}
\externaldocument[GE-]{01_law_of_graphic_equalisation_theory_referenced}
\externaldocument[PH-]{01a_graphic_equalisation_physics_grounding}
\externaldocument[BN-]{01d_graphic_equalisation_born_recovery}
\externaldocument[TP-]{01g_graphic_equalisation_total_perspective_vortex_grounding}
\externaldocument[NM-]{03_law_of_narrative_momentum_theory_referenced}

\newtheorem{definition}{Definition}[section]
\newtheorem{proposition}[definition]{Proposition}
\newtheorem{hypothesis}[definition]{Hypothesis}
\theoremstyle{remark}
\newtheorem{remark}[definition]{Remark}

\newcommand{\Rel}{\mathfrak R}
\newcommand{\Poss}{\mathcal P}
\newcommand{\Hist}{\mathcal H}
\newcommand{\Load}{\mathcal C}
\newcommand{\Obs}{\mathcal O}

\title{The Law of Graphic Equalisation --- Information Gravity Grounding:\\
\large Change-Load, Quantum Realisation, and the Reconciliation Programme}
\author{Matthew Parslow\\Independent Researcher}
\date{\textbf{Draft \texorpdfstring{\textemdash}{---}\ \PaperDraft}\\[2pt]\PaperDisplayDate \texorpdfstring{\textbullet}{\textperiodcentered}\ document revision \PaperRevision}

\usepackage{fancyhdr}
\pagestyle{fancy}
\fancyhf{}
\fancyfoot[C]{\thepage}
\fancyfoot[R]{\scriptsize\PaperFooterStamp}
\renewcommand{\headrulewidth}{0pt}
\fancypagestyle{plain}{\fancyhf{}\fancyfoot[C]{\thepage}%
  \fancyfoot[R]{\scriptsize\PaperFooterStamp}\renewcommand{\headrulewidth}{0pt}}
% GENERATED by tools/build_paper_metadata.py from paper_metadata.json.
% DO NOT EDIT: edit paper_metadata.json, then regenerate this file.
\providecommand{\PaperDraft}{draft-2026-09-06.1}
\providecommand{\PaperDisplayDate}{6 September 2026}
\expandafter\def\csname PaperRevision@00_is\endcsname{129}
\expandafter\def\csname PaperRevision@01_law_of_graphic_equalisation_theory_referenced\endcsname{294}
\expandafter\def\csname PaperRevision@01a_graphic_equalisation_physics_grounding\endcsname{47}
\expandafter\def\csname PaperRevision@01b_graphic_equalisation_mind_and_self_grounding\endcsname{47}
\expandafter\def\csname PaperRevision@01c_graphic_equalisation_cosmology_grounding\endcsname{11}
\expandafter\def\csname PaperRevision@01d_graphic_equalisation_born_recovery\endcsname{6}
\expandafter\def\csname PaperRevision@01e_graphic_equalisation_communication_grounding\endcsname{6}
\expandafter\def\csname PaperRevision@01f_graphic_equalisation_qualia_and_consciousness_grounding\endcsname{3}
\expandafter\def\csname PaperRevision@01g_graphic_equalisation_total_perspective_vortex_grounding\endcsname{4}
\expandafter\def\csname PaperRevision@01h_graphic_equalisation_information_gravity_quantum_reconciliation\endcsname{5}
\expandafter\def\csname PaperRevision@01i_graphic_equalisation_passive_continuation_momentum_gravity\endcsname{12}
\expandafter\def\csname PaperRevision@01s_graphic_equalisation_statement\endcsname{294}
\expandafter\def\csname PaperRevision@02_seldon_manifolds_theory_referenced\endcsname{31}
\expandafter\def\csname PaperRevision@02a_seldon_manifolds_psychological_grounding\endcsname{15}
\expandafter\def\csname PaperRevision@02b_seldon_manifolds_social_and_psychological_application\endcsname{7}
\expandafter\def\csname PaperRevision@02c_seldon_manifolds_decision_making\endcsname{11}
\expandafter\def\csname PaperRevision@03_law_of_narrative_momentum_theory_referenced\endcsname{32}
\expandafter\def\csname PaperRevision@04_narrativium_theory_referenced\endcsname{23}
\expandafter\def\csname PaperRevision@05_frames_relation_and_the_empty_limit\endcsname{11}
\expandafter\def\csname PaperRevision@06_assessing_an_account_of_this_kind\endcsname{29}
\expandafter\def\csname PaperRevision@07_causal_attitude_networks_on_the_substrate\endcsname{4}
\expandafter\def\csname PaperRevision@08_certified_answers_as_an_audit\endcsname{4}
\expandafter\def\csname PaperRevision@09_adaptability\endcsname{94}
\expandafter\def\csname PaperRevision@09a_adaptability_implementation\endcsname{7}
\expandafter\def\csname PaperRevision@10_applied_graphic_equalisation\endcsname{7}
\expandafter\def\csname PaperRevision@10a_red_queen_race_civilisational_rate_of_change_llms_biological_timescales_anthropocene_transition\endcsname{13}
\expandafter\def\csname PaperRevision@11_relational_information_ethics\endcsname{16}
\expandafter\def\csname PaperRevision@A1_constructive_graphic_equalisation_implementation_referenced\endcsname{8}
\expandafter\def\csname PaperGenerated@01s_graphic_equalisation_statement\endcsname{1}
\edef\PaperMetadataStem{\jobname}
\expandafter\ifx\csname PaperRevision@\PaperMetadataStem\endcsname\relax
  \PackageError{paper_metadata}{No metadata for \PaperMetadataStem}{Add the paper stem to paper_metadata.json.}
  \def\PaperRevision{?}
\else
  \expandafter\let\expandafter\PaperRevision\csname PaperRevision@\PaperMetadataStem\endcsname
\fi
\newcommand{\PaperDateLine}{\textbf{Draft \texorpdfstring{\textemdash}{---}\ \PaperDraft}\\[2pt]\PaperDisplayDate\ \texorpdfstring{\textbullet}{\textperiodcentered}\ document revision \PaperRevision}
\newcommand{\PaperStatementDateLine}{\textbf{Draft \texorpdfstring{\textemdash}{---}\ \PaperDraft}\\[2pt]\PaperDisplayDate\ \texorpdfstring{\textbullet}{\textperiodcentered}\ statement}
\expandafter\ifx\csname PaperGenerated@\PaperMetadataStem\endcsname\relax
  \newcommand{\PaperFooterStamp}{\textsc{\PaperDraft\,r\PaperRevision}}
\else
  \newcommand{\PaperFooterStamp}{\textsc{\PaperDraft\ / statement / from long r\PaperRevision}}
\fi
\newcommand{\PaperNoticeStamp}{\textbf{DRAFT --- \PaperDraft\ r\PaperRevision} (\PaperDisplayDate).}


\begin{document}
\maketitle
\begin{center}
\fbox{\parbox{0.92\textwidth}{\small
\PaperNoticeStamp\ This is the \emph{01h} grounding companion to \emph{The
Law of Graphic Equalisation --- Theory}, \emph{Physics Grounding},
\emph{Born Recovery}, \emph{Total Perspective Vortex Grounding}, and
\emph{The Law of Narrative Momentum}. It presents a falsifiable
\emph{reconciliation programme},
not a completed solution to quantum gravity. Its positive result is a typed
bridge ledger: it says what a physical realization must supply before the
project can claim a common quantum--gravitational domain. It derives no metric,
Einstein equation, Newton constant, Lorentz factor, numerical force law,
dark-matter phenomenology, or successful quantum-gravity theory. Physical
gravity is treated as one possible grounding of a more general
relational-information principle; self-grounding and cross-modal recurrences
remain independent structural witnesses with separate tests, not evidence for
physical gravity.}}
\end{center}
\vspace{1em}

\begin{abstract}
Graphic Equalisation supplies a relational vocabulary for possibility,
frame-indexed realization, history, and change cost. Its physics companion
proposes a common change-budget shape for motion and change-load; its Born note
maps the vocabulary to a conditional quantum-reconstruction programme. Neither
map supplies a physical state space, gauge-invariant record, conserved source,
metric, field equation, exact clock law, or empirical discriminator. This paper
makes that boundary operational. It defines a common relational realization
situation, separates inherited statements from new physical hypotheses, and
sets out the Q1--Q6 quantum and G1--G8 gravitational bridges a concrete model
must close. The result is a programme that can fail: a candidate loses its
reconciliation claim if it cannot provide a noncircular quantum composition and
Born bridge, a conserved dimensioned source, a controlled geometry and
proper-time recovery, compatibility with established limits, and a constrained
discriminator from general relativity plus standard quantum-field-theory
effective field theory. Related thermodynamic, holographic,
information-geometric, causal-set, and gravity-mediated-entanglement work
supplies important conditional results and targets; none supplies these missing
maps for the present substrate. The paper therefore advances a test protocol,
not an asserted unification.
\end{abstract}

\section{Scope and terminology translation}\label{sec:scope}
The familiar question is whether a relational-information account can reconcile
quantum descriptions with gravity. The answer established here is precise: the
current corpus supplies a \emph{programme specification}, not a quantum-gravity
solution. A reconciliation requires one model with translation maps and
compatible predictions in a shared quantum--gravitational domain. An analogy
shares a motif; a recovery instantiates a target formalism under declared
assumptions; a derivation proves a target from stated premises. These are not
interchangeable achievements.

The project starts from a general relational-information principle. Physical
gravity is a proposed grounding of that principle, not its source. This
asymmetry matters: a physical failure can reject a proposed physical grounding
without deciding every nonphysical structural use of the relational vocabulary.
Likewise, recurrence in self-model, cognitive, computational, or institutional
cases is a structural witness only when its own comparison class and tests have
been supplied. It neither proves a shared physical mechanism nor transports a
gravitational result across modalities.

\paragraph{Terminology guard.} ``Information'' below is not thermodynamic
entropy, von Neumann entropy, a modular Hamiltonian, energy, or stress--energy.
A frame-relative record is not automatically a diffeomorphism-invariant local
observable. A graph or partial order is not a Lorentzian manifold; a statistical
or Fisher metric is not the dynamical spacetime metric; and an entanglement
pattern is not metric dynamics. Each identification would require a stated map
and an independently checkable domain.

\section{Existing project premises and the common reconciliation object}\label{sec:object}
The inherited vocabulary can be collected without promoting it to physical
ontology. Let the \emph{relational realization situation} be
\[
 \Rel_t=(W,G_t,\Gamma_t,\Poss_t,\Hist_t,\Load_t,\Obs_t),
 \qquad G_t=(g_t,r_t,\Gamma_t),
\]
where $W$ is an individuating frame relation; $\Poss_t=\operatorname{Poss}(G_t)$
is the admissible-continuation family; $\Hist_t$ is realized history when it is
represented; $\Load_t$ is change-cost bookkeeping; and $\Obs_t$ is the relevant
frame-held record/readout relation. The inherited realization relation is
\[
 \operatorname{Coll}^{W}_t:\Poss_t\rightsquigarrow G_{t+1}.
\]
This is a typed common schema, not a state space, action, path integral,
causal-set dynamics, or observable algebra.

\begin{proposition}[Frame-indexed realization]\label{prop:relative-realization}
Under the corpus definitions, realization is indexed to the frame holding the
relevant interaction and record. No frame-independent global realization time
follows from this index.
\end{proposition}
\begin{proof}
The possibility family, admissibility grammar, record, and represented
successor are all defined relative to the antecedent graph as held through $W$.
Erasing $W$ erases the relation that individuates the update. This establishes a
typed internal consequence only; it does not establish covariant measurement,
no-signalling, complete positivity, gauge invariance, or an observer-independent
causal order.
\end{proof}

The object is useful because it prevents category drift. Possibility has a slot
in $\Poss_t$, change-load in $\Load_t$, local succession in realization order,
and physical geometry in a future map from $\Rel_t$. It is insufficient because
none of those slots fixes amplitudes, a conserved source, proper time, or a
metric. The history and cost slots inherit their limited roles from
\emph{The Law of Narrative Momentum}: operative history and alteration burden
are not universal physical energy. Their comparison-class, support, unit, and
ordering discipline inherits from \emph{Total Perspective Vortex Grounding}.
Neither companion supplies a physical source or observable; both prevent that
missing bridge from being hidden. The table records the paper's claim ledger
before any bridge is named.

\begin{center}\small
\begin{tabular}{@{}>{\raggedright\arraybackslash}p{0.22\textwidth}>{\raggedright\arraybackslash}p{0.23\textwidth}>{\raggedright\arraybackslash}p{0.24\textwidth}>{\raggedright\arraybackslash}p{0.23\textwidth}@{}}
\toprule
\textbf{Derived / inherited claims} & \textbf{New hypotheses} & \textbf{Open bridge obligations} & \textbf{Empirical predictions / constraints}\\
\midrule
A typed relational schema can be assembled; realization is frame-indexed; admissible possibility does not determine amplitudes. & H-QA: a physical amplitude/state representation exists for the schema. & Define composition, transformations, normalization, and gauge-invariant instruments. & The resulting probabilities must meet the stated quantum tests; a competing non-Born theory satisfying the premises rejects a claimed Born derivation.\\
\addlinespace
Change-cost is a grounding-dependent relation, not an identified physical source; cycle order is not proper time. & H-CL: change-load maps to a dimensioned conserved effective source. H-CT: one clock map yields both tested dilation limits. & Supply source units, exchange/conservation law, geometry, dynamics, and clock observable. & Weak-field, equivalence, local-Lorentz, redshift, Shapiro-delay, and precision-clock tests constrain or reject the map.\\
\addlinespace
Relational graphs do not determine Lorentzian geometry or dynamics; 01d remains conditional at its linearity/composition seam. & H-GEO: a controlled coarse-graining yields physical geometry and observables. & Recover a continuum/alternative, metric/connection or replacement, QFT/EFT overlap, and black-hole scope. & A candidate must state a finite parameterization and at least one result distinguishable from GR plus standard QFT EFT; a null or contrary result must exclude it.\\
\bottomrule
\end{tabular}
\end{center}

\section{Physics grounding: change-load and existing relativistic limits}\label{sec:physics}
The physics companion proposes a fixed change-budget picture: motion and
internal change-load are two uses of one throughput. That is an information-level
shape intended to organize a reading of kinematic and gravitational clock
behavior. It is not a physical source equation. In particular, the existing
project has not derived a metric, a connection, $G$, Einstein equations, a
Lorentz transformation group, the exact Lorentz factor, a gravitational-redshift
law, a numerical force law, or dark-matter phenomenology.

\begin{proposition}[No source identification follows]\label{prop:no-source}
The inherited alteration-cost and change-load vocabulary does not define an
energy density, momentum flux, stress tensor, modular energy, entropy, or source
for curvature.
\end{proposition}
\begin{proof}
The cost currency, payer, and comparison class vary by grounding. A physical
source requires units, an operational measurement, a locality or nonlocality
domain, a conservation or exchange law, and a coupling rule. None is fixed by
$\Load_t$ or by the budget image. Thus no insertion into a gravitational field
equation is typed by the present premises.
\end{proof}

\begin{proposition}[No clock law follows]\label{prop:no-clock-law}
Local realization order or closure-cycle counting does not entail either the
special-relativistic Lorentz factor or gravitational redshift.
\end{proposition}
\begin{proof}
An order admits many inequivalent clock functions. The present corpus supplies
no metric, worldline functional, transformation group, or physical clock
observable that identifies an operational counter with measured proper time. In
particular, a closure-count is only the default simple clock hypothesis: the
corresponding execution may be locally closure-driven, externally clocked, or
scheduled by a superframe without an internally available discriminator.
\end{proof}

\begin{hypothesis}[H-CL: load--source bridge]\label{hyp:load-source}
A candidate realization supplies a map from $\Load_t$ to a dimensioned effective
source $T_{ab}^{(\mathrm{eff})}$, or to an explicitly different named source,
with operational definition, coupling, and conservation or accounted exchange
law.
\end{hypothesis}
\begin{hypothesis}[H-CT: unified-clock bridge]\label{hyp:clock}
A candidate realization defines a physical clock observable and proves, in its
claimed regime, both the tested kinematic and gravitational clock relations, or
states an experimentally distinct alternative with its measurement protocol.
\end{hypothesis}

These hypotheses are hard constraints, not labels for a desired answer. H-CL
fails if no source can meet the required conservation compatibility, or if a
nonuniversal coupling violates established equivalence or local-Lorentz bounds.
H-CT fails if its physical clocks disagree with established dilation, redshift,
Shapiro-delay, post-Newtonian, or precision-comparison results in the regime it
claims. The empirical review literature documents how tightly metric-gravity and
universality departures are constrained \citep{will2014}.

\section{Quantum possibility, realization, and what 01d supplies}\label{sec:quantum}
The Born note supplies a correspondence programme between project vocabulary and
operational/reconstruction conditions. It identifies a classical probability
backbone and marks complex structure, composition, and linearity assumptions or
residuals. This paper carries that boundary forward rather than treating a
quantum vocabulary as a completed noncircular derivation.

\begin{proposition}[Possibility is not amplitude]\label{prop:no-amplitude}
From $G_t$, $\Gamma_t$, and $\Poss_t=\operatorname{Poss}(G_t)$ alone, no complex
amplitude representation follows.
\end{proposition}
\begin{proof}
Membership in an admissible set supplies neither a scalar field nor addition,
phase, inner product, normalization, composition, or dynamics. Inequivalent
probability theories can be defined over the same alternative set.
\end{proof}

\begin{proposition}[Conditional Born boundary]\label{prop:born-boundary}
If a realization supplies quantum state structure, a suitable noncontextual
probability measure, and every dimensional and composition assumption of an
applicable reconstruction result, then Born weights can follow within that
result's stated domain. The present corpus does not establish those premises in
full.
\end{proposition}

\begin{hypothesis}[H-QA: quantum representation bridge]\label{hyp:quantum}
A candidate realization provides an operational state/effect space, an amplitude
or explicitly alternative representation, sequential and composite composition,
and normalized frame-transformation rules without importing the conclusion it
claims to derive.
\end{hypothesis}

Quantum information tools cannot close this gap merely by being invoked.
Holographic codes, information geometry, and gravity-mediated-entanglement
protocols ordinarily use quantum state spaces and Born statistics as input. A
claimed Born bridge fails if it assumes Hilbert linearity, the trace rule, or
standard composition while announcing those structures as outputs. A physical
record bridge also fails if the record changes under gauge redundancy: gravity
makes local observables and subsystem factorization nontrivial
\citep{donnelly2016}.

\section{The joint bridge map and exact unclosed obligations}\label{sec:bridges}
A reconciliation candidate must publish a concrete realization $\mathfrak S$ of
$\Rel$, not a prose analogy. The following maps are the minimum interfaces:
\[
 \begin{aligned}
 Q_{W,t}&:\Poss_t\longrightarrow\mathcal S(\mathcal A_W), &
 A_{W,t}&:\Poss_t\longrightarrow\mathcal H_W,\\
 \Phi_\epsilon&:\Rel\longrightarrow(M,g_{ab},\nabla,\Obs_{\rm phys}), &
 \Theta_W&:\Rel\longrightarrow\tau_W .
 \end{aligned}
\]
Here $\mathcal S(\mathcal A_W)$, $\mathcal H_W$, $(M,g_{ab},\nabla)$, and
$\tau_W$ are targets to be defined by the candidate, not objects supplied by the
corpus. The amplitude codomain may be replaced by another operational theory,
but then the alternative must state its composition, probability, and test
domain.

\paragraph{Quantum obligations.} Q1 defines operational state/effect spaces and
normalized affine probabilities. Q2 defines sequential and composite
composition, including any gauge-constrained replacement for naive
factorization. Q3 derives or admits a complex Hilbert/C$^*$-algebraic structure
with transformations and composition. Q4 derives Born weights from Q1--Q3 plus
all necessary assumptions, or openly admits them as an input. Q5 makes records
and realization events gauge-invariant relational observables or instruments.
Q6 specifies unitary or nonunitary domain, normalization, and strong-gravity or
black-hole scope.

\paragraph{Gravitational obligations.} G1 supplies a controlled coarse-graining
to a Lorentzian continuum or stated alternative with causal observables. G2
identifies metric/connection data or an alternative and their transformation
law. G3 closes H-CL. G4 supplies a field equation or tested alternative with
error and domain. G5 derives or parameterizes universality and local Lorentz
behavior. G6 closes H-CT. G7 derives entropy/area and black-hole behavior or
limits the domain explicitly. G8 identifies how gauge-observable and relevant
no-go assumptions are satisfied or evaded.

\paragraph{Conditional compatibility target T0.} Let $\mathfrak S$ be a
concrete realization. If it closes Q1--Q6 and G1--G8, preserves the claimed
operational predictions in their shared domain, and publishes a discriminator
from GR plus standard QFT EFT, it is a candidate reconciliation in that domain.
T0 is a target specification rather than a theorem of this paper: no concrete
$\mathfrak S$ has been supplied, so none of its antecedents has been proved.

\section{Relation to prior art}\label{sec:prior-art}
\paragraph{Thermodynamic and entropic routes.} Jacobson derives the Einstein
equation as an equation of state under a local horizon Clausius relation,
area-entropy, and stress-energy-flux assumptions \citep{jacobson1995}. His
entanglement-equilibrium route adds fixed-volume causal-diamond and quantum
input \citep{jacobson2016}. These are conditional gravitational results, not a
derivation that $\Load_t$ is entropy, energy, or curvature. Verlinde's
entropic-force proposal is also a proposal with additional screen, temperature,
and equipartition assumptions \citep{verlinde2011}; implementations motivated by
later emergent-gravity work faced empirical mismatches in isolated dwarf systems
\citep{lelli2017}. This paper claims neither their force law nor their
phenomenology.

\paragraph{Holography and quantum error correction.} AdS/CFT is a conjectured
gauge/gravity duality with extensive evidence in particular settings
\citep{maldacena1998}. In static settings, Ryu--Takayanagi relates boundary
entropy to a minimal bulk surface under its own conditions \citep{ryu2006}; the
covariant Hubeny--Rangamani--Takayanagi extension uses extremal bulk surfaces
\citep{hubeny2007}. Entanglement-wedge reconstruction has quantum-error-correcting
structure in a code subspace \citep{almheiri2015}. These furnish exact examples
of encoding and reconstruction, not evidence that an arbitrary relational graph
produces generic spacetime, quantum amplitudes, or metric dynamics.

\paragraph{Information geometry and causal sets.} Information geometry gives
metrics on statistical models \citep{amari2016}; those metrics are not thereby
Lorentzian spacetime metrics. Causal-set theory gives a concrete locally finite
partial-order programme for spacetime \citep{bombelli1987,surya2019}. It is a
useful target because causal order and volume provide more physical structure
than the present graph vocabulary, but it does not itself supply this project's
Born, budget, or cross-modal bridges.

\paragraph{Quantum-gravity witnesses and constraints.} Gravity-mediated
entanglement proposals test nonclassical mediation only conditional on their
specified locality and mediator assumptions \citep{bose2017,belenchia2018}.
Even a successful witness would not select this substrate, detect a graviton, or
validate a unique quantum-gravity theory. Any emergent massless-spin-two claim
must also identify its relation to the assumptions of Weinberg--Witten
\citep{weinberg1980}; ``emergence'' alone is not an evasion. At low energies,
quantum general relativity is usefully organized as an effective field theory
\citep{donoghue1994}, which a candidate must recover or delimit.

\section{Empirical discriminators, recovery conditions, and falsifiers}\label{sec:falsifiers}
No numerical prediction is claimed because no realization fixes a source
normalization, continuum scale, parameter, or observable. A candidate earns a
calculation only after it gives a quantity with units, parameterization,
uncertainty/control analysis, and comparison to GR plus standard QFT EFT in the
same domain. Dimensional rhetoric is not a prediction.

\paragraph{Recovery conditions.} The classical recovery condition is a specified
low-velocity, weak-field, decohered limit with source law and force, potential,
or geodesic equation. The special-relativistic condition is covariant
observables and the exact tested clock relation. The general-relativistic
condition is a metric/effective dynamics with weak-field, redshift, light, and
post-Newtonian agreement. The quantum condition is a state space, composition,
dynamics, instruments, and Born statistics. The semiclassical condition is a
renormalized observable and controlled backreaction domain. None is presently
obtained by the corpus; each is a required future calculation.

\paragraph{Failure protocol.} The programme direction fails as a claimed
reconciliation for a candidate if any of the following occurs:
\begin{enumerate}[leftmargin=*]
\item no noncircular quantum-composition and Born bridge can be constructed;
\item no dimensioned, conserved source or accounted exchange law maps
change-load to physical dynamics;
\item no continuum/alternative map supplies gauge-invariant causal and clock
observables compatible with local-Lorentz and equivalence constraints;
\item the candidate conflicts with established clock, weak-field, or EFT
predictions within its stated regime;
\item it meets a no-go theorem's assumptions while asserting the prohibited
object; or
\item all proposed realizations remain observationally identical to GR plus
standard QFT EFT and provide no further finite, rejectable parameter claim.
\end{enumerate}
The final item is not a demand for novelty theatre. It distinguishes a compatible
reinterpretation from a theory with additional empirical content. A valid
proposal must name the observable, parameter region, comparator, protocol, and
null or contrary outcome that excludes its mechanism.

\section{Conclusion}\label{sec:conclusion}
The project currently supports one hard result: a reconciliation claim has a
checkable type. A physical candidate must translate frame-indexed realization
into gauge-invariant instruments, admissible possibility into a quantum
representation without circular Born input, change-load into a conserved
physical source, and relational structure into controlled causal geometry and
proper-time observables. It must then recover the domains where GR and quantum
field theory already work and risk a constrained disagreement.

Until those bridges are published for a concrete realization, this is not a
quantum-gravity solution. It is a reconciliation programme with named
obligations and failure conditions. Physical gravity can ground one realization
of the wider relational-information principle; it cannot be used as that
principle's source, and nonphysical structural witnesses must continue to earn
their status under their own tests.

\begin{thebibliography}{99}
\raggedright
\bibitem[Almheiri, Dong and Harlow(2015)]{almheiri2015}
A.~Almheiri, X.~Dong and D.~Harlow, ``Bulk locality and quantum error correction
in AdS/CFT,'' \emph{Journal of High Energy Physics} \textbf{04}, 163 (2015),
arXiv:1411.7041.

\bibitem[Amari(2016)]{amari2016}
S.-i.~Amari, \emph{Information Geometry and Its Applications} (Springer, 2016),
https://doi.org/10.1007/978-4-431-55978-8.

\bibitem[Belenchia et al.(2018)]{belenchia2018}
A.~Belenchia \emph{et al.}, ``Quantum superposition of massive objects and the
quantization of gravity,'' \emph{Physical Review D} \textbf{98}, 126009 (2018),
https://doi.org/10.1103/PhysRevD.98.126009.

\bibitem[Bombelli et al.(1987)]{bombelli1987}
L.~Bombelli, J.~Lee, D.~Meyer and R.~D.~Sorkin, ``Space-time as a causal set,''
\emph{Physical Review Letters} \textbf{59}, 521 (1987),
https://doi.org/10.1103/PhysRevLett.59.521.

\bibitem[Bose et al.(2017)]{bose2017}
S.~Bose \emph{et al.}, ``Spin entanglement witness for quantum gravity,''
\emph{Physical Review Letters} \textbf{119}, 240401 (2017),
https://doi.org/10.1103/PhysRevLett.119.240401.

\bibitem[Donnelly and Giddings(2016)]{donnelly2016}
W.~Donnelly and S.~B.~Giddings, ``Observables, gravitational dressing, and
obstructions to locality and subsystems,'' \emph{Physical Review D} \textbf{94},
104038 (2016), https://doi.org/10.1103/PhysRevD.94.104038.

\bibitem[Donoghue(1994)]{donoghue1994}
J.~F.~Donoghue, ``General relativity as an effective field theory,''
\emph{Physical Review D} \textbf{50}, 3874 (1994),
https://doi.org/10.1103/PhysRevD.50.3874.

\bibitem[Hubeny, Rangamani and Takayanagi(2007)]{hubeny2007}
V.~E.~Hubeny, M.~Rangamani and T.~Takayanagi, ``A covariant holographic
entanglement entropy proposal,'' \emph{Journal of High Energy Physics}
\textbf{07}, 062 (2007), arXiv:0705.0016.

\bibitem[Jacobson(1995)]{jacobson1995}
T.~Jacobson, ``Thermodynamics of spacetime: The Einstein equation of state,''
\emph{Physical Review Letters} \textbf{75}, 1260 (1995),
https://doi.org/10.1103/PhysRevLett.75.1260.

\bibitem[Jacobson(2016)]{jacobson2016}
T.~Jacobson, ``Entanglement equilibrium and the Einstein equation,''
\emph{Physical Review Letters} \textbf{116}, 201101 (2016),
https://doi.org/10.1103/PhysRevLett.116.201101.

\bibitem[Lelli et al.(2017)]{lelli2017}
F.~Lelli, S.~S.~McGaugh, J.~M.~Schombert and M.~S.~Pawlowski, ``Testing emergent
gravity with isolated dwarf galaxies,'' \emph{Monthly Notices of the Royal
Astronomical Society} \textbf{468}, L68 (2017), arXiv:1706.00785.

\bibitem[Maldacena(1998)]{maldacena1998}
J.~M.~Maldacena, ``The large N limit of superconformal field theories and
supergravity,'' \emph{Advances in Theoretical and Mathematical Physics}
\textbf{2}, 231 (1998), arXiv:hep-th/9711200.

\bibitem[Ryu and Takayanagi(2006)]{ryu2006}
S.~Ryu and T.~Takayanagi, ``Holographic derivation of entanglement entropy from
AdS/CFT,'' \emph{Physical Review Letters} \textbf{96}, 181602 (2006),
https://doi.org/10.1103/PhysRevLett.96.181602.

\bibitem[Surya(2019)]{surya2019}
S.~Surya, ``The causal set approach to quantum gravity,'' \emph{Living Reviews
in Relativity} \textbf{22}, 5 (2019),
https://doi.org/10.1007/s41114-019-0023-1.

\bibitem[Verlinde(2011)]{verlinde2011}
E.~Verlinde, ``On the origin of gravity and the laws of Newton,'' \emph{Journal
of High Energy Physics} \textbf{04}, 029 (2011), arXiv:1001.0785.

\bibitem[Weinberg and Witten(1980)]{weinberg1980}
S.~Weinberg and E.~Witten, ``Limits on massless particles,'' \emph{Physics
Letters B} \textbf{96}, 59 (1980),
https://doi.org/10.1016/0370-2693(80)90610-3.

\bibitem[Will(2014)]{will2014}
C.~M.~Will, ``The confrontation between general relativity and experiment,''
\emph{Living Reviews in Relativity} \textbf{17}, 4 (2014),
https://doi.org/10.12942/lrr-2014-4.
\end{thebibliography}

\end{document}
